\uuid{6BZg}
\exo7id{4140}
\titre{exo7 4140}
\auteur{quercia}
\organisation{exo7}
\datecreate{2010-03-11}
\isIndication{false}
\isCorrection{false}
\chapitre{Fonction de plusieurs variables}
\sousChapitre{Dérivée partielle}
\cpge{mpsi-mp2i.fonctions-deux-variables.b, pcsi.fonctions-deux-variables.b, ptsi.fonctions-deux-variables.b, mpsi-mp2i.fonctions-deux-variables.c, pcsi.fonctions-deux-variables.c, ptsi.fonctions-deux-variables.c, bcpst1.fonctions-deux-variables.c, bcpst1.derivees-primitives-integrales.b, mp-mpi.calcul-differentiel-optimisation.a, mp-mpi.calcul-differentiel-optimisation.b, pc.fonctions-plusieurs-variables.c, psi.fonctions-plusieurs-variables.c}
\module{Analyse}
\niveau{L2}
\difficulte{}

\contenu{
\question{
Soit $f : {\R^2} \to \R$ de classe $\mathcal{C}^1$ et
$g : {\R^n} \to \R, {(x_1,\dots,x_n)} \mapsto
{ \int_{t=0}^1 f(t,x_1t+x_2t^2 + \dots + x_nt^n)\,d t.}$

Montrer que $g$ est de classe $\mathcal{C}^1$ et calculer ses dérivées partielles.
}
}
